\documentclass[10pt]{article}

\usepackage[utf8]{inputenc}

\usepackage[T1]{fontenc}

\usepackage{amsmath}

\usepackage{amsfonts}

\usepackage{amssymb}

\usepackage[version=4]{mhchem}

\usepackage{stmaryrd}

\usepackage{graphicx}

\usepackage[export]{adjustbox}

\graphicspath{ {./images/} }

\usepackage{bbold}



\begin{document}

ма льным многочлен ом әлемента $\boldsymbol{\alpha}$ в. Расширение $K / k$ наз. а л г е 6 р а и ч е с к и м, если всякий әлемент из $K$ алгебраичен над $k$. Р., не являюЩееся алгебраическим, наз. т р а в с ц е н д е в т в Ы м. Р. наз. н о р м а л ь н ы м, если оно алгебраическое и всякий неприводимый в $k[x]$ многочлен, обладающий корнем в $K$, разлагается в $K[x]$ на линейные множители. Подполе $k$ наз. а л г е 6 р а и ч е с к и 3 а м к н ут ы м в $K$, если каждый алгебраический над $k$ әлемент из $K$ на самом деле лежит в $k$, т. е. всякий элемент из $K / k$ трансцендентен над $k$. Поле, алгебраически замкнутое в любом своем Р., является алаебраически замкнутыл полем.



Расширение $K / k$ наз. к о н е ч н о по $\mathrm{p} о$ ж д е нвы м (или расши рением ко неч но го т и п а), если в $K$ существует такое конечное подмножество элементов $S$, что $K$ совпадает с наименьшим подполем, содержащим $S$ и $k$. В этом случае говорят, что $K$ порождается множеством $S$ над $k$. Если $K$ может быть порождено над $k$ множеством из одного алемента $\alpha$, то Р. наз. П р о с т ы м и обозначается $K=k(\alpha)$. Простое алгебраич. расширение $k(\alpha)$ полностью определяется минимальным многочленом $f_{\alpha}$ порождающего әлемента $\alpha$. Точнее, если $k(\beta)-$ другое простое алгебраическое Р. и $f_{\alpha}=f_{\beta}$, то существует изоморфизм расширений $k(\alpha) \rightarrow k(\beta)$, переводящий $\alpha$ в $\beta$. Далее, для любого неприводимого многочлена $f \in k[x]$ существует простое алгебраич. расширение $k(\alpha)$ с минимальным многочленом $f_{\alpha}=f$. Оно может быть построено как факторкольцо $k[x] / f k[x]$. С другой стороны, для всякого простого трансцендентного расширения $k(\alpha)$ существует изоморфизм расширений $k(\alpha) \rightarrow k(x)$, где $k(x)$ - поле рациональных функций от $x$ над $k$. Любое Р. конечного типа может быть получено с помощью конечной цепочки простых Р.



Расширение $K / k$ наз. к о н е ч н ы м, если $K$ как алгебра конечномерна над полем $k$, и б е с к о н е чн ы м, если эта алгебра бесконечномерна. Размерность этой алгебры наз. с т е п е н ь ю р с ш и р е н и я $K / k$ п обозначается $[K: k]$. Каждое конечное Р. является алгебраическим и каждое алгебраическое Р. конечного типа - конечным. Степень простого алгебраического Р. совпадает со степенью соответствующего минимального многочлена. Напротив, простое трансцендентное $\mathrm{P}$. бесконечно.



Пусть дана последовательность расширений $K \subset L \subset$ $\subset M$. Расширение $M / K$ является алгебраическим тогда и только тогда, когда и $L / K$ и $M / L$ - алгебраические Р. Далее, $M / K$ конечно тогда и только тогда, когда конечны $L / K$ и $M / L$, причем



\[

[M: K]=[M: L] \cdot[L: K]

\]



Если $P / k$ и $Q / k$ - два алгебраических P. и $P Q$ композит полеїі $P$ и $Q$ в нек-ром их общем надполе, то $P Q / k$ также алгебраическое $\mathrm{P}$.



См. также Сепарабельное расширение, Трансцендентное расширение.



лит.: [1] Б у р б а к и Н., Алгебра. Многочлены и поля. Упорядоченные группы, пер. с франц., М., 1965; [2] В а н д е p В а p д е н Б. JI., Алгебра, 2 изд., пер. с нем., М., 1979; [3] т. 1, пер. с англ., М., 1963; [4] Л е н г С., Алгебра, пер. с англ., т. ', 1968. с англ., М., 1963, [4] л е н г С., Алгебра, Пер. с англ.,



\includegraphics[max width=\textwidth, center]{2023_07_08_74e965fc06e8ebfdf1abg-1}

с т р а н с т в а $X$ - топологическое пространство $Y$, в к-ром $X$ является всюду плотным подпространством. Если $Y$ бикомпактно, то оно наз. би к о м п а к тны м р а с ши и е н и м, если $Y$ хаусдорфовох у у со рфо вым р а сши р ением.



РАСІІИРЕНИЯ ОБЛ.И. Войцеховский. РАСІІИРЕНИЯ ОБЛАСТИ ПРИНЦИП, п р и нц и п К а р л е м а н а: гармоническая мера $\omega(z, \alpha, D)$ дуг $\propto$ границы $\Gamma$ области $D$ может только возрастать при распирении области $D$ через дополнительные дуги $\beta \subset \Gamma, \alpha \bigcup \beta=\Gamma$. Точнее, пусть граница $\Gamma$ области $D$ на плоскости комплексного переменного $z$ состоит из конечного числа жордановых кривых, $\alpha$ - часть $\Gamma$, состоящая из конечного числа дуг $\Gamma$, и пусть область $D^{\prime}$ есть р а с ши р ен и е о б л а с т и $D$ через дополнительные дуги $\beta=\Gamma \backslash \alpha$, то есть $D \subset D^{\prime}$ и $\alpha$ есть часть границы $\Gamma^{\prime}$ области $D^{\prime}$. Тогда для гармонич. мер справедливо неравенство $\omega(z, \alpha, D) \leqslant \omega\left(z, \alpha, D^{\prime}\right), \quad z \in D$, причөм знак равенства здесь имеет место только в случае $D^{\prime}=D$. P. о. п. справедлив и для гармонич меры относительно областей евклидова пространства $\mathbb{R}^{n}, n \geqslant 2$, или $\mathbb{C}^{n}, n \geqslant 1$.



Р. о. П. находит важные применения в различных ситуациях, связанных с оценками гармонич. меры. Напр., еще Т. Карлеман [1] дал при помощи Р. о. П. решение



\includegraphics[max width=\textwidth, center]{2023_07_08_74e965fc06e8ebfdf1abg-1(1)}

Пусть граница $\Gamma$ односвязноӥ области $D$ состоит из конечного числа жордановых дуг, точка $\zeta$ расположена на $\Gamma$ или $\zeta \in C \bar{D}, \Delta=\{z:|z-\zeta|<R\}-$ круг радиуса $R$ с центром $\zeta$, а $\alpha$ - часть $\Gamma$, попавшая в $\Delta_{R}=\Delta \cap D$. Требуется найти оценку снизу для гармонич. меры $\omega\left(z, \alpha, \Delta_{R}\right)$, зависящую только от $R$ и $|z-\zeta|, z \in \Delta_{R}$. Решение выражается неравенством



\[

\omega\left(z, \alpha, \Delta_{R}\right) \geqslant \frac{2}{\pi} \operatorname{arctg}\left(\frac{2}{\theta(R)} \ln \frac{R}{|z-\zeta|}\right),

\]



где $R \theta(R)$ - сумма длин дуг пересечения



\[

\{z:|z-\zeta|=R\} \cap D .

\]



Поскольку $\theta(R) \leqslant 2 \pi$, то



\[

\omega\left(z, \alpha, \Delta_{R}\right) \geqslant \frac{2}{\pi} \operatorname{arctg}\left(\frac{1}{\pi} \ln \frac{R}{|z-\zeta|}\right), z \in \Delta_{R} .

\]



Имеются обобщения проблемы Карлемана - Мийю п уточнения формул (1), (2) (см. [3]). Р. о. і. позволяет доказать также Линделёба теоремы. Многочисленные применения Р. о. п. и формул типа (1), (2) дал А. Мийю (см. [2], а также [3], [4]).



Jum.: [1] C a r l e m a n T., "Ark. Mat. Astron. Fys.", 1921, v. 15, N $\mathcal{N}_{10}$ 1 [2] M i 111 o u x H., "J. math. pures et appl.», 9 sér., 1924, ․ $3 ;$; [3] Н е в а н ли и н н а Р., Однозначные аналитичеекие функции, пер. с нем., М.- Л., $1941 ;$ [4] Е в г р а ф о в М. А., Аналитические функции, 2 изд., М., 1968. Е Д Соломенцев.



РАСШИРЕННАЯ КОМПЛЕКСНАЯ ПЛОСКОСТЬ плоскость комплексного переменного $\mathbb{C}$, компактифицированная посредством добавления бесконечно удаленной точки $\infty$ и обозначаемая $\overline{\mathbb{C}}$. Окрестностью $\infty$ является внешность любого круга в $\overline{\mathbb{C}}$, т. е. множество вида $\{\infty\} \cup\left\{z \in \mathbb{C}:\left|z-z_{0}\right|>R\right\}, \quad R \geqslant 0$. Р. к. п. есть Александрова бикомпактное расширение плоскости $\mathbb{C}$, гомеоморфное и конформно эквивалснтное Римана cøepe. С ф е риче с к а я, или хо р д а льн а я, ме т р и к а на $\overline{\mathbb{C}}$ дается формулами



\[

\begin{gathered}

\rho(z, w)=\frac{2|z-w|}{\sqrt{1+|z|^{2}} \sqrt{1+|w|^{2}}}, z, w \in \mathbb{C}, \\

\rho(z, \infty)=\frac{2}{\sqrt{1+|z|^{2}}} .

\end{gathered}

\]



Лит.: [1] М а р к у ші е в и ч А. И., Теория аналитических функций, 2 изд., т. 1-2, М., $1967-68 ;$ [2] ШI а б а т Б. В., Введение в комплексный анализ, 2 изд., ч. $1 \frac{2}{E}$, М. М., 1976.



РАСІЕПЛЕНИЯ МЕТОД - сеточный метод решения нестационарных задач со многими пространственными переменными, в к-ром переход от заданного временного слоя $t_{n}$ к новому слою $t_{n+1} t_{n}+\tau$ осуществляется за счет последовательного решения сеточных аналогов родственных нестационарных задач с меньшим числом пространственных переменных (см. [1] - [4]).





\end{document}